\documentclass{article}
\usepackage{amsmath,amssymb,amsthm}
\newtheorem{theorem}{Theorem}
\newtheorem{lemma}[theorem]{Lemma}
\newtheorem{corollary}[theorem]{Corollary}
\newcounter{paperpremise}
\newcounter{definingequation}
\newenvironment{proofstep}{\par\smallskip\ifdim\leftskip<3em\relax\advance\leftskip by 1em\relax\fi\noindent\ignorespaces}{\par\smallskip}
\setlength{\textwidth}{6.5in}
\setlength{\oddsidemargin}{0in}
\setlength{\emergencystretch}{\linewidth}
\title{Two arithmetic bounds}
\author{}
\date{}
\begin{document}
\maketitle
\section{Results}
\begin{theorem}[{Weighted bound}]\label{theorem:1}
Let \(x\) and \(y\) be rational numbers. Assume \(2 \cdot x + y \le 9\). Assume \(1 \le y\). Then \(x \le 4\).
\end{theorem}

\begin{proof}
Suppose for contradiction that \(x > 4\).
Adding \(2 \cdot x + y - 9 \le 0\), \(1 - y \le 0\) and \(2 \cdot \left(4 - x\right) < 0\) gives \(2 \cdot x + y - 9 + \left(1 - y\right) + 2 \cdot \left(4 - x\right) < 0\).
The terms on the left cancel, so the inequality would give \(0 < 0\), a contradiction.


\end{proof}

\begin{theorem}[{Rounded bound}]\label{theorem:2}
Let \(x\) and \(y\) be integers. Assume \(2 \cdot x \le y\). Assume \(y \le 9\). Then \(x \le 4\).
\end{theorem}

\begin{proof}
Suppose for contradiction that \(\neg \left(x \le 4\right)\).
Taking the linear combination of \(y - 2\,x \ge 0\) and \(y \le 9\) with respective coefficients 1 and -1 gives \(-2\,x \ge -9\).
Divide \(-2\,x \ge -9\) by \(-2\), reversing the order. Since \(x\) is integral, \(x \le \left\lfloor\frac{9}{2}\right\rfloor = 4\).
From \(\neg \left(x \le 4\right)\) we get \(4 < x\), hence \(x \ge 5\) over the integers.
The two bounds require \(5 \le x \le 4\).
This is impossible because \(5 > 4\).


\end{proof}

\end{document}
